%%%%%
%
%% Dans ce fichier, les �ventuels accents sont cod�s
%% en ISO-8859-15 (aka latin9)
%% La ligne suivante donne trois exemples 
%% e accent aigu : �, e accent grave : �, e-dans-l'o :
\section{Incidence � partir du vent horizontal et vertical}

\itemstyle{Op\'erateur~:} \operfmt{Cal\_angl\_inc}
\index{Cal\_angl\_inc}

%%%%%OBJET%%%%
\itemstyle{Objet~:}
Obtenir le vecteur angle d'incidence mesur� entre l'�coulement
 et le m�t pour ensuite corriger de la distortion.

%%%%%ARGUMENTS%%%%
\bigskip
%% Dans ce fichier, les �ventuels accents sont cod�s
%% en ISO-8859-15 (aka latin9)
%% La ligne suivante donne trois exemples 
%% e accent aigu : �, e accent grave : �, e-dans-l'o : �
%%%%%ENTREES%%%%
\begin{tabarg}
\multicolumn{3}{l}{\titretab {Entr\'ees}}\\
\hline
$UU$ & Vecteur, \texttt{float} & Vent moyen dans l'axe moyen (unit� indif�rente)\\
$WW$ & Vecteur, \texttt{float} & Composante w verticale (m�me  unit� que $UU$)\\
\end{tabarg}

\medskip
%%%%%SORTIES%%%%
\begin{tabarg}
\multicolumn{3}{l}{\titretab {Sorties}}\\
\hline
$i_m$ & Vecteur, \texttt{float} & Angle d'incidence mesur� (degres)\\
\end{tabarg}


%%%%%FORMULE%%%%
\itemstyle{Formule}

%% Formule de l'operateur : Cal_angl_inc
%% Dans ce fichier, les �ventuels accents sont cod�s
%% en ISO-8859-15 (aka latin9)
%% La ligne suivante donne trois exemples 
%% e accent aigu : �, e accent grave : �, e-dans-l'o : �
\[
i_m = 
\begin{cases}
\arctan(\frac{WW}{UU}) \qquad \text{ si } UU \ne 0 \\
\text{valeur manquante} \qquad \text{ sinon }  \\
\end{cases}
\]



%%%%%COMMENTAIRE%%%%
\itemstyle{Commentaire~:}
Cet op�rateur est consid�r� comme obsol�te : 
\begin{itemize}
\item le nom est trop sp�cifique
\item il ne permet pas de choisir l'unit� du parametre calcul�
\end{itemize}
On utilisera plut�t l'op�rateur \hyperref[OP-Atan2]{\operfmt{Atan2}}

%%%%%AUTEUR%%%%
\itemstyle{Auteur~:}
CNRM/GMEI/TRAMM

%%%%%REFERENCE%%%%
%\itemstyle{R\'ef\'erences~:}
%INCONNU
